% !TEX root = ../monograph.tex
\subsection{The stationary integral with its cutoff retained}
\label{ch09:stationary-remainder}

We now give the local analytic implication behind the auxiliary bound.
It is an application of Laplace's method, for which we refer to
\cite[Chapter II]{Wong2001}, with three features that must be checked
here: the integration dimension grows with the auxiliary period, the
Hessian at the critical configuration depends on the coarse field, and the integration
cutoff cannot be discarded before estimating its complement.
All actions and matrices in the following argument belong to the same
\emph{pure} small-field block integral, the block integral of the path action alone: incoming activities are omitted, the
prescribed path action and block kernel are unchanged, and the fine-field
indicator $\mathbf1_{\{\chi(U_p)\le t_f\ \text{for every fine }p\}}$
is retained, where $t_f=t_{m+1}$. Thus its integrand is zero on the fine large-field
complement. The other sectors at the fine scale have already been attributed to the
large-field expansion.

Put $z=\widehat\beta_{m+1}$, $g=\beta'^{-1/2}$ and
$h=z^{-7/16}$. The letter $z$ is a scalar inverse coupling parameter in this
subsection; the Gaussian integration variable is $\xi$. Auxiliary
periods and coordinate counts are bounded by fixed powers of
$1+\log z$. Write $M=M_{m+1}^{\rm cut}$ for the cutoff amplitude
in \eqref{ch04:thresholds}, and take $D=d_0M$ with a fixed small
$d_0>0$. We shall use the following two distinct facts:
\[
 M=z^{o(1)},\qquad (1+\log z)^k=o(M^\epsilon)
 \quad\text{for every fixed }k,\epsilon>0.
\]
Both follow directly by taking logarithms of
$M=\exp((\log\log z)^2)$.

Here and below $z^{o(1)}$ in an upper bound means that for every
$\epsilon>0$ the quantity is $O_\epsilon(z^\epsilon)$ as $z\to\infty$
at fixed stage and block ratio.

Let $x$ be the coarse logarithmic coordinates, with their supremum
norm, and $y\in\RR^n$ the free fine coordinates after fixing the block axial
gauge. The classical action including the penalty term of the block kernel is
$\phi(x,y)$; write $\rho(x)$ for its critical configuration and
$A(x)=D_y^2\phi(x,\rho(x))$ for its positive Hessian.
The extra exponent $\psi(x,y)$ includes the fine loop term and minus
the logarithm of the product Haar-chart Jacobian. Suppose the critical-configuration map $\rho$,
action, extra exponent and relative logarithmic determinant have the
fixed derivative orders required below on a common neighborhood,
with bounds $z^{o(1)}$. The Hessian gap is uniform in the auxiliary
volume at the fixed block ratio. These are the stationary and path
estimates of Chapters~\ref{ch05:gaussian}--\ref{ch07:matching}; no
different flat covariance is substituted here.

For a fine-coordinate displacement $g\xi$, with $\xi\in\RR^n$, set
$s(x)=\tfrac12\log\det A(x)$ and define the second-order Taylor remainder
divided by $g^2$ through its integral formula
\[
 V(g,x,\xi)=\int_0^1(1-t)
 D_y^2\phi(x,\rho(x)+tg\xi)[\xi,\xi]\,dt.
\]
Stationarity gives
$\phi(x,\rho(x)+g\xi)-\phi(x,\rho(x))=g^2V(g,x,\xi)$,
including the limiting value
$V(0,x,\xi)=\tfrac12\xi^*A(x)\xi$.
With $\gamma_0=N(0,A(0)^{-1})$, put
\begin{align}
 a(g,x,\xi)&=s(x)-s(0)+V(0,0,\xi)-V(g,x,\xi)\notag\\
 &\quad+\psi(x,\rho(x))-\psi(x,\rho(x)+g\xi),\notag\\
 R(g,x)&=\int_{|\xi|\le D}e^{a(g,x,\xi)}\,d\gamma_0(\xi),
 \qquad L(g,x)=\log R(g,x).
 \label{ch09:truncated-integral}
\end{align}
The determinant factor is essential: at $g=0$ the exact identity is
\[
 R(0,x)=N(0,A(x)^{-1})\{\,|\xi|\le D\,\},
\]
which is the Gaussian measure of a ball, not identically one. Changing $\xi$ to
$-\xi$ also gives the exact parity $R(-g,x)=R(g,x)$.

Here are the uniform bounds needed for this elementary application
of Laplace's method. Differentiating the divided difference on
$|\xi|\le D$ costs a fixed polynomial in $D$ and the stationary
derivative bounds. It is therefore $z^{o(1)}$ at every fixed order.
The exponent vanishes at $(g,x)=(0,0)$, so on
$|g|,\|x\|_\infty\le h$ it is at most one in absolute value for
large $z$. If the reference Gaussian measure of the ball is at least $1/2$, then
$e^{-1}/2\le R\le e$, and the same derivative bounds $z^{o(1)}$ hold
for $L$. To obtain derivatives through order $r$ and the two
derivatives in $g$ used for parity, it suffices to retain derivatives
of $\phi,\rho$ through order $r+4$ and of the other terms through
order $r+2$. These are finite-stage requirements.

If $A(x)^{-1}\le cI$ and $n$ is at most a fixed polynomial in
$\log z$, a coordinate Gaussian tail bound gives
\[
 1-R(0,x)\le 2n\exp\bigl(-D^2/(2cn)\bigr).
\]
This is smaller than every fixed inverse power of $z$.
The same conclusion for its fixed derivatives follows either by
the Gaussian score formula or by finite-order interpolation on a
ball of radius $z^{-1}$ inside the parameter domain. We recall the
needed interpolation fact: a superpolynomially small value and
polynomially bounded derivatives through order $r+1$ imply
superpolynomially small derivatives through order $r$ on a smaller
ball. It follows successively from Taylor's estimate
$\|DF\|\le2\sup|F|/u+u\sup\|D^2F\|$ on a radius-$u$ ball,
applied to successive derivatives with shrinking radii.
No coordinate-dimension factor is needed when operator norms are
used. Parity and Taylor's theorem in $g$ now yield, for every fixed
$K>0$,
\begin{equation}
 \|D_x^r L(g,x)\|\le C_r(z)\beta'^{-1}+z^{-K},
 \qquad C_r(z)=z^{o(1)}.
 \label{ch09:even-gain}
\end{equation}
This includes $r=0$; the cutoff error at $g=0$ has been retained
throughout the argument.

\subsection{Identifying the retained region in the pure block integral}
\label{ch09:stationary-exterior}

Write $\mathcal D_x(y)$ for the configuration map of the block axial gauge.
Consider the region
$\mathcal S_x=\mathcal D_x(\rho(x)+g\overline B(0,D))$.
Let $\mathscr Z(x)$ be the pure block integral, and set
\[
 \Phi(x)=\beta'\phi(x,\rho(x))+\psi(x,\rho(x))
                         +\tfrac12\log\det A(x).
\]
If $N$ is the number of free fine links and $c_H$ the Haar density
constant, its normalization constant is
$\Omega=\pi_{\beta'}(\sqrt{2\pi}g)^n c_H^N$;
$\pi_{\beta'}$ is the block-kernel normalizing prefactor.
When $\mathcal S_x$ lies in the chart, all its fine plaquettes satisfy
the regularity condition $\chi(U_p)\le t_f$, and the actions for all values of the auxiliary index with positive probability weight coincide with $\phi$, the
change of variables gives exactly
\begin{equation}
 \int_{\mathcal S_x}\text{(pure integrand)}
                  =\Omega e^{-\Phi(x)}R(g,x).
 \label{ch09:retained-identity}
\end{equation}
The probability weights of the auxiliary index sum to one, so there is no additional factor counting its values.
The hypotheses about this region must be established before
using the identity.

Here is a verification using the geometry and the Hessian gap.
By the estimate \eqref{ch05:nonlinear-tree} in the block axial gauge, a fine
configuration lies in a fixed closed coordinate ball of the chart\mgubi{Issue 10: in which norm? existence/ takes the supremum ball and uses its convexity.}, the \emph{safe chart}, whenever its fine
plaquettes and the deviations of the coarse links from the values of the averaging operation are below a fixed $\delta>0$.
On that chart assume $D_y^2\phi\ge\gamma I/2$. Stationarity gives
\[
 \phi(x,y)-\phi(x,\rho(x))
                      \ge\gamma|y-\rho(x)|^2/4.
\]
Thus a chart point outside $\mathcal S_x$ adds at least
$\Delta=\gamma D^2/4$ to the exponent. A point outside the chart
with all fine plaquettes small must have such a deviation exceeding
$\delta$, and contributes a fixed positive multiple of $\beta'\delta^2$.
The possibly negative cubic, quartic and loop terms are bounded by
the path coefficient bounds. If $W$ bounds the auxiliary coordinate and site
counts, their total contribution to the exponent is at most
\[
 W B_{\rm loop},\qquad
 B_{\rm loop}=\beta'C(t_f^3+t_f^4)+\Lambda t_f^2,
 \qquad 0\le\Lambda\le2M.
\]
All such terms and the logarithmic Haar/determinant normalization
term are $o(M^2)$. In particular, the normalization term
$-N\log[c_H(\sqrt{2\pi}g)^{\dim\gLie}]$ is only a polynomial
in $\log z$. The bound on $\phi(x,\rho(x))$ on $\|x\|<h$
also leaves the fixed penalty for a deviation exceeding $\delta$ available, because
that value tends to zero and $\Delta/\beta'\to0$.
Consequently the exterior integrand is bounded by
$\Omega e^{-\Phi(x)}e^{-\Delta+o(M^2)}$.

It remains to verify the regularity condition on the fine plaquettes on the retained ball. For the
auxiliary configurations used here, which are reflected periodic extensions (pullbacks of the original field to an auxiliary torus under the coordinatewise reflection map $\pi_v$ of Section~\ref{ch09:local-density}), this can be done directly
from the induced map on curvatures, without a lower bound from a sector
expansion. A full coordinate cycle of a reflected periodic extension traverses an
original path and then its inverse. Its holonomy is therefore exactly
one at every transverse base point. Equivalently, choose a coordinate
axial gauge on the original cube and pull that gauge back under the
reflection map. Plaquette filling bounds give logarithmic coordinates
with
\[
 \|x\|_\infty\le P(B)t_c,\qquad
 \|d_{1,c}x\|_\infty\le C t_c+P(B)t_c^2,
 \qquad t_c=t_m.
\]
The second estimate follows by expanding each plaquette product in
these small link coordinates. Crucially, its leading constant is
independent of $B$; the polynomial loss occurs only in the quadratic
error. This argument would fail for an arbitrary small-curvature
periodic field with nontrivial winding holonomies.

Let $\mathcal U(x)=\mathcal D_x(\rho(x))$ be the critical configuration as a physical
fine field. Its linearization is the representative $W_F$ in the block axial gauge of
Chapter~\ref{ch05:gaussian}. Its curl agrees with that of the smooth
representative $W_a$. The curvature identity and the
fine-row bound of the induced map $T$ on curvatures give
\[
 d_{1,f}W_a=H_f^{-1}TH_c d_{1,c},\qquad
 \|H_f^{-1}TH_c\|_{\infty\to\infty}\le C_jL^{-2}.
\]
The operator norm here sums each fine-output row over coarse
inputs; the opposite column sum has a different scaling.
The quadratic remainder of the critical-configuration map and the plaquette
product is bounded using their common second derivatives. Thus
\begin{equation}
 \sup_p\chi(\mathcal U(x)_p)
       \le C_jL^{-2}t_c+P_{j,L}(B)t_c^2.
 \label{ch09:stationary-cutoff-margin}
\end{equation}
Choose $L$ so that $C_jL^{-2}t_c\le t_f/8$, using the uniform
bound $t_c/t_f\le2$. Then increase the lower bound on $z$ so that
the second term is at most $t_f/8$. These choices are in the
required order: the constant in the first term is fixed before
$L$, while the polynomial in the second term may depend on $L$.

Plaquette functions have a fixed Lipschitz constant $L_0$ in the
physical fine-link supremum coordinates. Local chart maps have
fixed-stage derivative bounds. Consequently variation of $x$ by
$z^{-1}$ and of the fine coordinates at the critical configuration by $g\xi$,
$|\xi|\le D$, changes every plaquette distance from the identity by at most
$C_{j,L}z^{-1}+L_0gD$, after incorporating the fixed chart norm
conversion into $L_0$. Choose $d_0$ sufficiently small, so that
$L_0gD\le t_f/8$, and then raise the lower bound on $z$ for the term
$C_{j,L}z^{-1}$. Every point of each retained region now satisfies the regularity condition on the fine plaquettes,
with a strict margin. The common chart and the ball on which the actions for all values of the auxiliary index with positive probability weight coincide
contain these regions because $h+gD\to0$.
We have therefore obtained, on a full microscopic neighborhood of
each tested coarse configuration satisfying the regularity condition,
\begin{equation}
 0\le\mathscr Z(y)-\Omega e^{-\Phi(y)}R(g,y)
       \le\Omega e^{-\Phi(y)}e^{-c_0M^2},\qquad c_0>0.
 \label{ch09:exterior-bound}
\end{equation}
This deduction uses the linearization of the critical configuration (through its smooth representative $W_a$), its second
derivatives, the compatible path bounds and the geometric
estimate \eqref{ch05:nonlinear-tree}. The direct proof of the regularity condition applies to the
auxiliary fields that are reflected periodic extensions and to their microscopic neighborhoods, which
are precisely the configurations used for localization.

\subsection{Centering the full logarithm}
\label{ch09:centered-log}

Put $E=\log\mathscr Z-\log\Omega+\Phi-L(g,\cdot)$.
Since $R\ge e^{-1}/2$, \eqref{ch09:exterior-bound} gives
$0\le E\le2e\,e^{-c_0M^2}$. To pass to derivatives, write
the positive pure integrand at a fixed fine field and a fixed value of the auxiliary index as
$a(\omega)e^{p(x,\omega)}$, with $a\ge0$ independent of $x$.
Only the coarse-link kernel is differentiated. Its exponent has
fixed-order derivatives bounded by $C_r\beta' n_b$, where
$n_b$ is the number of coarse links. The exponential chain rule
and integration give relative bounds
\[
 \|D^r\mathscr Z(x)\|\le P_r(\beta'n_b)\mathscr Z(x).
\]
The logarithm chain rule thus gives polynomial bounds without
dividing an absolute error by a possibly tiny partition function.
The other terms in $E$ have polynomial derivative bounds from the
stationary estimates. Interpolation on the microscopic neighborhoods
in \eqref{ch09:exterior-bound} proves that every fixed jet of $E$
is smaller than every inverse power of $z$.

Let $a_c=\mathcal K$ be the coarse classical action and
$d_c=\mathcal L+\alpha\mathcal Q$ the prescribed coarse quantum contribution:
the inherited loop action of the coarse path action and its quadratic part
multiplied by the coefficient $\alpha$ fixed by the matching, as in
\eqref{ch09:early}.\mgubi{Identified from the Lean definition YM.\allowbreak Precision.\allowbreak draft ($\alpha$ times the quadratic term plus the loop term), existence/\allowbreak stationary-output-construction.\allowbreak tex and v2 D-section; no $\Delta$ and no constant per site. Open: Lean uses $\mathcal K$, $\mathcal L$, $\mathcal Q$ without event indicators, while v3 defines them with indicators; they coincide on the small-field ball.} Define
\[
 C=a_c-\phi(\cdot,\rho),\qquad
 W=d_c-\psi(\cdot,\rho)-\tfrac12\log\det A.
\]
The corrected logarithm has the exact decomposition
\begin{equation}
 \log\mathscr Z+\beta'a_c+d_c
                  =\log\Omega+L+E+\beta'C+W.
 \label{ch09:pure-decomposition}
\end{equation}
Classical matching gives $D^rC(0)=0$ for $0\le r\le4$.
The quadratic quantum matching bounds $D^2W(0)$ by an
exponentially small finite-period error, with polynomial coordinate
norm costs. Its higher fixed derivatives are $z^{o(1)}$.
The full corrected logarithm has zero first derivative at zero
by constant gauge invariance and simplicity of $\gLie$.
Thus $DW(0)=-DL(0)-DE(0)$; no separate parity claim about each
determinant contribution is needed.

For a function on the radius-$h$ ball write its scaled derivative
bound through order $r$ as
$\max_{0\le k\le r}h^k\sup\|D^kf\|$.
Taylor's theorem, \eqref{ch09:even-gain}, and the two matching
cancellations show that the left side of
\eqref{ch09:pure-decomposition} minus its value at zero has this bound at most
\begin{equation}
 C_r(z)\bigl(\beta'^{-1}h+\beta'h^5+h^3\bigr)+z^{-K},
 \qquad C_r(z)=z^{o(1)},
 \label{ch09:pure-small-bound}
\end{equation}
where $K$ is arbitrary and the exponentially small quadratic
finite-period error has been included in the last term.
For $h=z^{-7/16}$ the three powers are respectively
$z^{-23/16}$, $z^{-19/16}$ and $z^{-21/16}$.
Their largest exponent is $-19/16<-23/20$.
Original small-field logarithmic coordinates and the directions admitted by the activity
seminorm \eqref{ch04:local-seminorm}, restricted to the cube $N_{\bar w}$ about the localization site and carried to the reflected periodic extension ($\bar w$ and the reflected extension are those of Section~\ref{ch09:local-density} below; the cube $N_r(x)$ is used there without a definition), are bounded by a polynomial in $B$ times
$t_m$, which is $o(h)$.\mgubi{All directions of the seminorm are meant, not only the affine ones (v2, D-section). v2 adds that the reflected torus must admit the charts and derivative bounds used here, and says that this and the finite-period comparison on a common analytic strip ``remain under verification''; v3 does not prove it either.} They therefore obey the same estimate.
This proves the auxiliary remainder needed for localization from
the explicit stationary, matching and actual-sector inputs above.
It does not require a formal expansion of an untruncated
exponential of a cubic Gaussian polynomial.
